\documentclass[uplatex,dvipdfmx]{jsarticle}

\usepackage[uplatex,deluxe]{otf} % UTF
\usepackage[noalphabet]{pxchfon} % must be after otf package
\usepackage{stix2} %欧文＆数式フォント
\usepackage[fleqn,tbtags]{mathtools} % 数式関連 (w/ amsmath)
\usepackage{hira-stix} % ヒラギノフォント＆STIX2 フォント代替定義（Warning回避）

\begin{document}

\title{加速度検知によるオーケストラシステムの構築}
\author{25G1021　梅澤聡颯太}
\date{2026年4月22日}
\maketitle
\section{技術的・社会的な背景}

現代において，現実世界の物理的な事象をセンサ等で取り込みデジタル処理を行う物理コンピューティングは，
スマートフォンやウェアラブルデバイスの普及により極めて身近なものとなった．加速度センサを用いて人間の動作を
検知する技術は一般化したが，その多くは「歩数計」や「画面の回転」といった，特定の数値を抽出するだけの定型的な
処理に留まっている．特にデジタル音楽やDTMの分野では，依然としてマウスによるクリックや
数値入力が主流であり，人間が持っている連続的な動きや呼吸に合わせた速度の変化といった
微細な表現をシステムに反映させることは，技術の進展に反して進んでいない．

\subsection{既存のシステム}

既存技術として任天堂が開発した「Wii Music」が存在する．加速度センサを搭載したWiiリモコンを楽器や指揮棒に見立てて
操作する，というコンセプトのゲームである\cite{ref:wii}．本システムでは，リモコンを上下に振る，あるいはボタンを押しながら傾ける
といった身体的な動作を加速度センサで検出し，その動きに合わせて仮想的な楽器の演奏や指揮を行いうことで
画面上のオーケストラの演奏テンポや音の強弱をリアルタイムに操作することが可能である．このシステムは，高度な演奏技術を習得していない者であっても，
直感的な身体動作のみで音楽表現を楽しむことができることを目的としたゲームである．主に家庭内でのエンターテインメントや，
リズム感・音楽性を養うための知育ツールとして，身体的インタフェースとデジタル音源を統合した事例と言える．

\subsection{自身のアイデア}

5台のArduinoを使用し，1台を指揮者，残り4台を演奏者として役割を分散させたオーケストラシステムを構築する．
指揮者役のArduinoには加速度センサを接続し，指揮棒の動きをリアルタイムで解析する．単に振った回数を
数えるのではなく，重力加速度を除いた動的な加速度ベクトルを算出することで，指揮の振り下ろしの瞬間を特定する．
この振り下ろし間の時間間隔を演算し，楽曲のBPMとして反映させる．さらに，加速度の大きさを音量の減衰率や音調のパラメータに
変換することで，指揮者の感情的な動きを音楽の強弱として表現する．


\begin{thebibliography}{9}

\bibitem{ref:wii}  任天堂. "Wii Music." https://www.nintendo.co.jp/wii/r64j/index.html, 2008.

\end{thebibliography}

\end{document} 
